一根长度为 \(L>0\) 的均匀细杆沿轴线发生小幅扭转，扭转刚度 \(C=GJ_p>0\)，单位长度转动惯量 \(\mu=\rho J_p>0\)。\(x=0\) 端固支，\(x=L\) 端刚接转动惯量 \(I\geq0\) 的转子，并通过刚度 \(K\geq0\) 的扭转弹簧接地。忽略耗散、剪切变形及轴向运动，令扭角为 \(\theta(x,t)\)，且至少 \(I,K\) 之一可非零。
(1) 从杆的角动量平衡和端部转子的力矩平衡推导波动方程及两端边界条件，写出严格的总能量守恒式。
(2) 用 \(u(x)e^{-i\omega t}\) 分离变量，推导非标准超越本征方程，定义 \(z=kL\)、\(p=KL/C\)、\(q=I/(\mu L)\)，并讨论正本征频率的离散性。
(3) 证明不同本征频率的模态在含端点转动惯量的加权内积下正交；推导含弹簧边界能的 Rayleigh 商，说明忽略端点惯量会造成什么错误。
(4) 求 \(I=K=0\) 时的全部本征频率；再令 \(I/(\mu L)\to\infty\) 而 \(K,C,L,\mu\) 固定，求最低频率的首项及其等效弹簧物理解释，并检验量纲。
